\documentclass[aspectratio=169,11pt]{beamer}

\usetheme{Madrid}
\usecolortheme{default}
\usefonttheme{professionalfonts}
\setbeamertemplate{navigation symbols}{}
\setbeamertemplate{footline}[frame number]

\usepackage{amsmath, amssymb, booktabs, array}

\title{Synthesis, Frontier Questions, and Student Research Designs}
\subtitle{Labor Markets and Political/Cultural Institutions -- Week 10}
\author{Teaching deck draft}
\date{}

\begin{document}

\begin{frame}
  \titlepage
\end{frame}

\begin{frame}{Central question}
\begin{itemize}
    \item What does a credible research design on institutions and labor markets look like?
    \item Why is historical, cultural, and institutional labor research worth doing now?
    \item Week 10 is a capstone studio: turn course themes into project architecture.
    \item The deliverable is a research memo, not a reading-list recap.
\end{itemize}
\end{frame}

\begin{frame}{Course architecture}
\small
\begin{tabular}{p{0.22\linewidth} p{0.30\linewidth} p{0.34\linewidth}}
\toprule
Course block & Institutional object & Labor margin \\
\midrule
Weeks 1--2 & law, contracts, enforcement, state capacity & rights, compliance, contract form \\
Weeks 3--4 & informality, voice, bargaining & formality, rents, wages, worker power \\
Weeks 5--6 & norms, culture, hierarchy, identity & search, sorting, mobility, access \\
Weeks 7--9 & persistence, reform, comparison, global systems & inequality, adjustment, portability \\
Week 10 & research design & mechanism, counterfactual, welfare \\
\bottomrule
\end{tabular}
\end{frame}

\begin{frame}{One institutional labor map}
\[
Y_i^{labor} = h(F_i, I_i, C_i, H_i, R_i, G_i; \theta)
\]

\small
\begin{tabular}{p{0.13\linewidth} p{0.70\linewidth}}
\toprule
Object & Interpretation \\
\midrule
$F$ & formal rules: law, courts, enforcement, bargaining, administration \\
$I$ & informal rules: norms, family systems, networks, expectations \\
$C$ & coercion, hierarchy, and durable social categories \\
$H$ & historical persistence and path dependence \\
$R$ & reform, implementation, compliance, adaptation \\
$G$ & comparative and global institutional environment \\
\bottomrule
\end{tabular}
\end{frame}

\begin{frame}{Why institutions, culture, and history matter now}
\begin{itemize}
    \item Labor markets are governed environments, not institution-free exchange.
    \item Formal institutions shape hiring, firing, payroll, contract form, bargaining, and protection.
    \item Informal institutions shape search, care, migration, occupation, trust, and acceptable work.
    \item Historical settings reveal mechanisms that contemporary cross sections often hide.
    \item Institutional detail is policy relevant when implementation, compliance, or beliefs mediate effects.
\end{itemize}
\end{frame}

\begin{frame}{Mechanism-first relevance}
\[
\text{Relevance} =
\text{Mechanism}+\text{Margin}+\text{Leverage}+\text{Portability}+\text{Welfare}
\]

\begin{itemize}
    \item Mechanism: what institutional channel changes behavior?
    \item Margin: what labor outcome is being explained?
    \item Leverage: why does the setting make the mechanism visible?
    \item Portability: what should travel beyond the setting?
    \item Welfare: who gains, who loses, and on what dimension?
\end{itemize}
\end{frame}

\begin{frame}{Defending a setting-specific paper}
\small
\begin{tabular}{p{0.27\linewidth} p{0.58\linewidth}}
\toprule
Skeptical question & Strong answer \\
\midrule
Why this historical episode? & It isolates a labor mechanism still relevant today. \\
Why one country? & The setting gives leverage on a broader labor question. \\
Why culture? & Norms shape labor supply, search, mobility, care, and occupation choice. \\
Why institutional detail? & Implementation, compliance, and enforcement are part of the treatment. \\
Why not just political economy? & The contribution is about work, wages, bargaining, allocation, and welfare. \\
\bottomrule
\end{tabular}
\end{frame}

\begin{frame}{Research architecture}
\[
\text{Contribution} =
\text{Question} + \text{Mechanism} + \text{Counterfactual}
+ \text{Data} + \text{Design} + \text{Interpretation}
\]
\begin{itemize}
    \item Question: what labor outcome is puzzling or important?
    \item Mechanism: what institutional channel generates the effect?
    \item Counterfactual: compared with what?
    \item Data and design: what source and variation make the claim credible?
    \item Interpretation: what travels, and what welfare object matters?
\end{itemize}
\end{frame}

\begin{frame}{Research-design template}
\scriptsize
\begin{tabular}{p{0.19\linewidth} p{0.27\linewidth} p{0.34\linewidth}}
\toprule
Element & Failure & Persuasive version \\
\midrule
Labor outcome & broad or loosely labor-related & maps to work, wages, bargaining, mobility, protection \\
Mechanism & institutional detail without channel & one institutional channel foregrounded \\
Counterfactual & no credible comparison & comparison matches the mechanism \\
Data & interesting but misaligned & observes labor margin and institution \\
Design & story outruns evidence & variation is explicit and defended \\
Portability & novelty substitutes for relevance & mechanism-first boundary conditions \\
Welfare & average effect only & distributional object is explicit \\
\bottomrule
\end{tabular}
\end{frame}

\begin{frame}{Design families}
\small
\begin{tabular}{p{0.34\linewidth} p{0.44\linewidth}}
\toprule
Design family & Typical institutional setting \\
\midrule
Legal or geographic discontinuity & boundaries, court rules, institutional assignment \\
Policy reform or staggered adoption & dismissal law, enforcement, transparency, bargaining \\
Enforcement or information shock & inspections, legal knowledge, court access \\
Historical administrative data & census, courts, maps, schools, payrolls, archives \\
Cross-region or cross-country comparison & legal regimes, informality, welfare states \\
Network or shift-share exposure & ancestry, migration, supply chains, sector composition \\
Linked employer-worker data & mobility, wages, compliance, incidence \\
\bottomrule
\end{tabular}
\end{frame}

\begin{frame}{Main critiques and honest responses}
\small
\begin{tabular}{p{0.28\linewidth} p{0.53\linewidth}}
\toprule
Critique & Response \\
\midrule
Bundled institutions & state what is identified and what remains bundled \\
Implementation differs & treat implementation as part of the object \\
Archives are selective & explain record construction and missing groups \\
Setting is narrow & map it to a broader labor taxonomy \\
Results may not travel & state boundary conditions \\
Persistence is endogenous & show why the labor channel remains plausible \\
Weak labor link & tie the setting to employment, wages, mobility, bargaining, welfare \\
\bottomrule
\end{tabular}
\end{frame}

\begin{frame}{Portability}
\[
\tau_s = \phi(M, Z_s, \Gamma_s)
\]
\begin{itemize}
    \item $M$: the mechanism, such as enforcement, coercion, networks, bargaining, migration access.
    \item $Z_s$: local institutional environment, such as coverage, state capacity, hierarchy, law.
    \item $\Gamma_s$: economic environment, such as market structure, informality, outside options.
    \item Mechanisms often travel more easily than point estimates.
    \item Broad relevance means knowing why something should or should not travel.
\end{itemize}
\end{frame}

\begin{frame}{What travels from anchor papers?}
\small
\begin{tabular}{p{0.25\linewidth} p{0.28\linewidth} p{0.31\linewidth}}
\toprule
Anchor & Mechanism that travels & What may not travel \\
\midrule
Dell & coercive labor institutions can persist through public goods and market access & Peruvian magnitude or exact channel mix \\
Naidu--Yuchtman & contract enforcement can reduce worker mobility and bargaining power & nineteenth-century legal and labor-market context \\
Munshi--Rosenzweig & networks can insure and constrain labor allocation & caste-specific institutional form \\
Beerli--Peri--Ruffner--Siegenthaler & migration rules alter matching and firm labor supply & Swiss geography and skill mix \\
Distelhorst--Hainmueller--Locke & private governance can affect supplier compliance & buyer commitment and audit credibility \\
\bottomrule
\end{tabular}
\end{frame}

\begin{frame}{Frontier directions}
\begin{itemize}
    \item Linked historical-to-modern census and administrative data.
    \item Enforcement and implementation measured with process data.
    \item Informality, contract form, payroll, and worker transitions.
    \item Culture, norms, networks, and social hierarchy inside labor markets.
    \item Comparative microdata and cross-border migration-regime designs.
    \item Global supply-chain governance, audits, certification, and worker voice.
    \item New forms of bargaining and representation beyond classic unions.
\end{itemize}
\end{frame}

\begin{frame}{Project memo template}
\small
\begin{tabular}{p{0.27\linewidth} p{0.55\linewidth}}
\toprule
Memo component & What it must say \\
\midrule
Labor outcome & work, wages, mobility, sorting, bargaining, protection, welfare \\
Institutional mechanism & rule, norm, hierarchy, enforcement, network, reform, governance \\
Unit of analysis & worker, firm, place, group, court, network, supply chain \\
Comparison group & untreated, less exposed, boundary neighbor, pre-reform, other regime \\
Identification strategy & variation, estimand, and primary threat \\
Data source & why the data see the mechanism and labor margin \\
Portability claim & what travels and what is local \\
Welfare object & incidence, inequality, autonomy, risk, protection, rents \\
\bottomrule
\end{tabular}
\end{frame}

\begin{frame}{Week 10 lab}
\begin{itemize}
    \item Reproduce: reverse-engineer Dell's historical institutional design.
    \item Diagnose: use Naidu--Yuchtman to separate coercive enforcement, mobility, and bargaining power.
    \item Transfer: adapt the architecture to a new institutions-and-labor project.
    \item Optional extension: use the Swiss migration opening as a modern portability case.
    \item Output: anchor menu, diagnostic table, and project memo scaffold.
\end{itemize}
\end{frame}

\begin{frame}{Closing bridge}
\begin{itemize}
    \item Dissertation-quality work usually starts with one disciplined mechanism.
    \item Institutions, culture, and history belong in labor economics when they explain labor allocation and welfare.
    \item A narrow setting can be powerful when it provides leverage on a broad labor question.
    \item The next step is to turn one setting into one mechanism and one mechanism into a research design.
\end{itemize}
\end{frame}

\begin{frame}{Takeaways}
\begin{itemize}
    \item Start with a labor outcome.
    \item Name the institutional mechanism.
    \item Define the counterfactual and design.
    \item Discuss portability honestly.
    \item State the welfare or distributional object.
    \item Let institutional detail sharpen the labor question.
\end{itemize}
\end{frame}

\end{document}
